% BHU PYQ -- Manually transcribed & error-corrected
% Paper: MTB-203A : Ancillary Mathematics I
% Exam: B.Sc. (Hons.) Semester II Examination, 2023-24

\documentclass[11pt,a4paper]{article}
\usepackage[a4paper,top=2.5cm,bottom=2.5cm,left=2.8cm,right=2.8cm,headheight=14pt]{geometry}
\usepackage{amsmath,amssymb}
\usepackage[T1]{fontenc}\usepackage[utf8]{inputenc}
\usepackage{lmodern,microtype,fancyhdr,enumitem,hyperref,lastpage,parskip}

\hypersetup{colorlinks=true,urlcolor=black,linkcolor=black,
  pdftitle={MTB-203A Ancillary Maths I -- BHU B.Sc. Sem II 2023-24}}
\pagestyle{fancy}\fancyhf{}
\renewcommand{\headrulewidth}{0.4pt}\renewcommand{\footrulewidth}{0.4pt}
\fancyhead[L]{\small   \textit{Banaras Hindu University}}
\fancyhead[R]{\small   \textit{MTB-203A: Ancillary Mathematics I}}
\fancyfoot[C]{\small Page  \thepage\ of \pageref{LastPage}}


\newcommand{\pts}[1]{\hfill   \textit{[#1]}}


\begin{document}

\begin{center}
  {\large\bfseries BANARAS HINDU UNIVERSITY}\\[2pt]
  {
\normalsize B.Sc.\ (Hons.) --- Semester II Examination, 2023-24}\\[6pt]
  \rule{0.8\linewidth}{0.5pt}\\[6pt]
  {\large\bfseries Paper No.\ MTB-203A: Ancillary Mathematics I}\\[4pt]
  {
\normalsize Time: 3 Hours \hfill Maximum Marks: 70}\\[2pt]
  {\small Ref.\ No.:\ 074894}
\end{center}
\rule{\linewidth}{0.4pt}
\smallskip



\noindent   \textit{Instructions: Answer   \textbf{five} questions including Question~1 which is compulsory. Figures in right-hand margin indicate marks.}
\bigskip




\noindent   \textbf{Question 1.}   \textit{(Compulsory --- 2 marks each.)}

\begin{enumerate}[label=(\alph*),leftmargin=*,itemsep=5pt,topsep=3pt]
  \item Draw the graph of $f(x)=|x|$ for $x\in\mathbb{R}$.
  \item If $(x+1,y-2)=(3,1)$, find $x$ and $y$.
  \item Let $A=\{1,2,3\}$, $B=\{3,4\}$, $C=\{4,5,6\}$. Find $A\cup(B\cap C)$.
  \item Let $A=\{1,2\}$ and $B=\{3,4\}$. Find the number of relations from $A$ to $B$.
  \item Let $V=\{a,e,i,o,u\}$ and $B=\{a,i,k,u\}$. Find $V-B$ and $B-V$.
  \item Show that $A\cup B=A\cap B$ implies $A=B$.
  \item Define diagonal matrix and give an example.
  \item For the $3 imes3$ matrix $A$ with rows $(1,2,3)$, $(1,3,2)$, $(2,4,7)$, find the minor $M_{12}$.
  \item Find $a,b,c,d$ from the equation $
\begin{pmatrix}2a+b & a-2b\\5c-d & 4d+3b\end{pmatrix}=
\begin{pmatrix}4&-3\\11&24\end{pmatrix}$.
  \item Given matrices $A$ with rows $(3,2,3)$, $(1,-5,3)$, $(2,2,-2)$ and $B$ with rows $(6,7,-5)$, $(-2,2,-6)$, $(0,1,7)$, find $A+B$.
  \item Find $AB$ where $A=
\begin{pmatrix}1\\-3\end{pmatrix}$ and $B=
\begin{pmatrix}2&-6&3\end{pmatrix}$.
\end{enumerate}
\medskip\rule{\linewidth}{0.2pt}

\medskip


\noindent   \textbf{Question 2.}

\begin{enumerate}[label=(\alph*),leftmargin=*,itemsep=5pt]
  \item Show that the relation $R=\{(m,n):|n-m| \text{ is even}\}$ on $X=\{3,4,5,6\}$ is an equivalence relation.\pts{6}
  \item Find the adjoint of $A=
\begin{pmatrix}1&2&3\\2&-5&1\\1&2&3\end{pmatrix}$.\pts{6}
\end{enumerate}
\medskip\rule{\linewidth}{0.2pt}

\medskip


\noindent   \textbf{Question 3.}

\begin{enumerate}[label=(\alph*),leftmargin=*,itemsep=5pt]
  \item If $A=
\begin{pmatrix}1&3&3\\3&-2&1\\4&1&-1\end{pmatrix}$, show $A^3-23A-40I=O$.\pts{6}
  \item Determine the row-rank of $A=
\begin{pmatrix}1&2\\2&3\\1&1\end{pmatrix}$.\pts{6}
\end{enumerate}
\medskip\rule{\linewidth}{0.2pt}

\medskip


\noindent   \textbf{Question 4.}

\begin{enumerate}[label=(\alph*),leftmargin=*,itemsep=5pt]
  \item Find the rank of $A=
\begin{pmatrix}1&2&-1&4\\2&4&3&6\\-1&-2&6&-7\end{pmatrix}$.\pts{6}
  \item Solve the system $x+y+z=6$, $y+3z=11$, $x-2y+z=0$.\pts{6}
\end{enumerate}
\medskip\rule{\linewidth}{0.2pt}

\medskip


\noindent   \textbf{Question 5.}

\begin{enumerate}[label=(\alph*),leftmargin=*,itemsep=5pt]
  \item Find the inverse of $A=
\begin{pmatrix}2&-1&4\\1&2&4\\5&1&2\end{pmatrix}$.\pts{6}
  \item Examine the consistency of $x+2y=2$, $2x+4y=3$.\pts{6}
\end{enumerate}
\medskip\rule{\linewidth}{0.2pt}

\medskip


\noindent   \textbf{Question 6.}

\begin{enumerate}[label=(\alph*),leftmargin=*,itemsep=5pt]
  \item Find minors and cofactors of all elements of the determinant
  $\Delta=
\begin{vmatrix}2&-3&5\\6&0&4\\1&5&-7\end{vmatrix}$.\pts{6}
  \item Show that $R=\{(5,5),(5,3),(2,2),(2,4),(3,5),(3,3),(4,2),(4,4)\}$ on $A=\{2,3,4,5\}$ is an equivalence relation.\pts{6}
\end{enumerate}
\medskip\rule{\linewidth}{0.2pt}

\medskip


\noindent   \textbf{Question 7.}

\begin{enumerate}[label=(\alph*),leftmargin=*,itemsep=5pt]
  \item Find the inverse of $A=
\begin{pmatrix}2&-1&3\\-1&3&-2\\2&-4&0\end{pmatrix}$ by row operations.\pts{6}
  \item Evaluate $\Delta=
\begin{vmatrix}1&2&4\\-1&3&0\\4&1&0\end{vmatrix}$.\pts{6}
\end{enumerate}

\vfill\rule{\linewidth}{0.4pt}

\begin{center}\small
\href{https://bhu-exam.vercel.app/}{Click here to visit our website}\\[4pt]
\href{https://me-aryan.vercel.app/}{Click here to visit our contributor}\\[4pt]
\href{https://quashan-me.vercel.app/}{Click here to visit our contributor}
\end{center}
\end{document}
